\documentclass[11pt]{article}

\usepackage{amsmath,amssymb,amsthm}
\usepackage{physics}
\usepackage{graphicx}
\usepackage{hyperref}
\usepackage{geometry}
\usepackage{booktabs}
\usepackage{enumitem}
\usepackage{bm}
\usepackage{float}

\geometry{margin=1in}
\graphicspath{{./figures/}}

% Theorem-like environments
\newtheorem{definition}{Definition}
\newtheorem{proposition}{Proposition}
\newtheorem{constraint}{Constraint}
\newtheorem{result}{Result}

\title{\textbf{The Hilbert Substrate Framework III:}\\
Complete Emergence of Spacetime and Matter\\
from Quantum Information Dynamics}

\author{Benjamin August Bray\\
\small Independent Researcher}

\date{January 2026}

\begin{document}

\maketitle

\begin{abstract}
We complete the Hilbert Substrate Framework by demonstrating the full emergence of three-dimensional space, fermionic matter, and local interactions from a geometry-free quantum substrate governed only by unitary dynamics and information-theoretic constraints. Building on Paper I (persistent heterogeneity in information propagation) and Paper II (accessibility phase transitions selecting spatial locality), we establish three central results.

First, we show that fermionic statistics emerge dynamically from a qubit substrate without any antisymmetrization postulate: Jordan--Wigner operators satisfy anticommutation relations to numerical precision, many-body spectra exhibit Fermi sea additivity, and ground-state energies display convexity characteristic of Pauli exclusion. Second, we identify a structural constraint---No-Refolding---that becomes unavoidable once stable matter exists: geometric transitions that would destroy fermionic structure are dynamically inaccessible under local unitary evolution, even when such transitions would reduce signaling cost. Third, through numerical experiments on locality recovery across interaction graphs of varying intrinsic dimension, we demonstrate that three-dimensional geometry provides the deepest accessible basin compatible with stable fermionic matter.

The substrate does not optimize over dimensions; it settles at the first geometry capable of supporting persistent structure. Once matter forms, No-Refolding locks the geometry in place. Three-dimensional space therefore emerges not as a fine-tuned input or anthropic selection, but as the first self-anchoring geometric configuration reachable under constrained unitary dynamics. Physical reality is a compromise: not the globally optimal description of the universe, but the simplest description accessible to systems containing stable matter.
\end{abstract}

\tableofcontents

\newpage

%==============================================================================
\section{Introduction}
%==============================================================================

\subsection{The Problem of Physical Structure}

The standard formulation of physics presupposes several structural features that are not derived from more fundamental principles: a three-dimensional spatial manifold, fermionic matter obeying the Pauli exclusion principle, and local interactions propagating at finite speed. While these features are well-confirmed empirically, their joint appearance raises foundational questions. Why does space have three dimensions rather than two, four, or some other number? Why do fermions exist at all, and what determines their statistics? Why are interactions local rather than nonlocal?

Traditional approaches treat these features as inputs---background assumptions or symmetry postulates that constrain the form of physical laws. Anthropic arguments note that stable atoms and chemistry require three dimensions with fermionic electrons, but do not explain how such a configuration is selected dynamically. String-theoretic approaches invoke extra compact dimensions, but presuppose a higher-dimensional geometric background. The problem of explaining why physical structure takes its observed form remains open.

\subsection{The Hilbert Substrate Framework}

The Hilbert Substrate Framework addresses these questions from a radically different direction. Rather than assuming spacetime, particles, or locality, it treats the universal Hilbert space and its unitary dynamics as fundamental, asking which forms of structure are \emph{unavoidable} under minimal information-theoretic constraints.

The guiding principle is conservative: new physical entities are not introduced unless required by the dynamics. Subsystems, geometry, particles, and forces are treated as emergent structures that arise from constraints on admissible operations within the Hilbert space. The framework seeks operational diagnostics capable of identifying when and how such structure becomes necessary.

This approach inverts the usual explanatory direction. Instead of asking ``How do quantum systems behave in a fixed spacetime background?'' we ask ``What spacetime and matter structures are dynamically selected from the space of all possible Hilbert space factorizations?''

\subsection{Summary of the Three-Paper Program}

This work completes a three-part investigation:

\textbf{Paper I: Persistent Heterogeneity in Information Propagation} \cite{HSF1}. We introduced the diagnostic of Heterogeneity in Information Propagation (HIP), which quantifies spatial variation in how localized perturbations influence neighboring subsystems under unitary evolution. Using numerical simulations on finite interaction graphs, we demonstrated that HIP generically develops nonuniform, temporally persistent structure. Certain regions of the substrate exhibit stable dominance in information transport, with rank ordering preserved across extended time intervals. This establishes that persistent structure in information flow can arise dynamically from minimal constraints, providing an operational order parameter for emergent organization.

\textbf{Paper II: Accessibility and the Emergence of Spatial Locality} \cite{HSF2}. We demonstrated the existence of an accessibility phase transition in the structure of unitary orbits. For small systems ($N \leq 3$), optimization flows reliably locate the global minimum (the ``Harmonion'' eigenbasis). However, beyond a critical system size ($N \geq 5$), this global minimum becomes kinetically inaccessible, and the system is trapped in a robust local basin corresponding to spatial geometry. We mapped the phase diagram of this transition and concluded that spatial locality emerges not because it is the simplest description of physics, but because it is the simplest \emph{accessible} description for large-scale systems.

\textbf{This paper} addresses the remaining questions: What determines which geometric structure is selected? Why does dimensionality freeze at three? And how does fermionic matter fit into this emergence chain? We demonstrate that fermionic statistics arise dynamically without antisymmetrization postulates, identify the No-Refolding constraint that locks geometry once matter exists, and show that three dimensions are uniquely selected as the first geometry where stable matter can form and anchor.

%==============================================================================
\section{Theoretical Framework}
%==============================================================================

\subsection{The Substrate Model}

We consider a finite-dimensional Hilbert space $\mathcal{H} = (\mathbb{C}^2)^{\otimes N}$ representing $N$ qubits. No assumption is made regarding an underlying spatial geometry or classical background; all structure arises from the pattern of interactions and the constraints imposed on admissible dynamics.

Locality of interactions is encoded by an undirected interaction graph $G = (V, E)$. Each vertex corresponds to a qubit subsystem, and an edge between two vertices indicates the presence of a direct coupling term in the Hamiltonian. The graph serves purely as an interaction topology and is not assumed to carry metric, dimensional, or geometric significance \emph{a priori}.

The system evolves under a Hamiltonian consisting of local terms and pairwise interaction terms:
\begin{equation}
H = \sum_{i \in V} H_i + \sum_{(i,j) \in E} H_{ij},
\end{equation}
where $H_i$ acts on a single subsystem and $H_{ij}$ couples subsystems connected by an edge.

\subsection{Information-Theoretic Constraints}

We work within a minimal framework defined by three information-theoretic constraints. These are not dynamical laws but structural conditions required for observer-accessible physics.

\begin{constraint}[No-Signaling]
Local operations on one subsystem cannot be used to transmit information instantaneously to spacelike-separated subsystems. Operationally, this bounds the rate at which correlations can be established across the interaction graph.
\end{constraint}

\begin{constraint}[No-Forgetting]
Global evolution is unitary; information is never destroyed, only redistributed within the total Hilbert space. Records may delocalize through entanglement, but they must persist within the global state.
\end{constraint}

\begin{constraint}[No-Refolding]
Once stable, structured matter exists, the accessible geometric configurations of the substrate are dynamically restricted. Existing structure cannot be arbitrarily refactored into alternative geometric arrangements without destroying the matter that anchors it.
\end{constraint}

The first two constraints are standard in quantum information theory. The third---No-Refolding---is not imposed independently but emerges as a consequence of applying No-Signaling and No-Forgetting to substrates containing persistent structure, as we demonstrate in Section \ref{sec:constraint_separation}.

\subsection{Locality Cost and the Unitary Orbit}

From the perspective of pure unitary dynamics, the specific tensor factorization of $\mathcal{H}$ is not fixed. Two Hamiltonians $H$ and $H' = UHU^\dagger$ are related by a unitary transformation and share identical spectra. However, they differ in their structure relative to a fixed local basis.

To quantify structural complexity, we expand the Hamiltonian in the Pauli basis:
\begin{equation}
H = \sum_{k=1}^{4^N} c_k P_k, \quad c_k = \frac{1}{2^N} \text{Tr}(H P_k),
\end{equation}
where $\{P_k\}$ is the $N$-qubit Pauli basis. We assign a locality weight $w(P_k)$ to each basis element, equal to its Hamming weight (the number of non-identity factors). The locality cost functional is:
\begin{equation}
\mathcal{C}_p(H) = \frac{\sum_k w(P_k)^p |c_k|^2}{\sum_k |c_k|^2}.
\end{equation}

Minimizing $\mathcal{C}_p(H)$ over the unitary orbit $\mathcal{O}_H = \{UHU^\dagger | U \in SU(2^N)\}$ is equivalent to finding the basis in which the Hamiltonian looks ``most local.''

%==============================================================================
\section{Emergent Fermionic Statistics}
%==============================================================================

\subsection{The ``Mush versus Steel'' Problem}

A universe composed solely of bosonic degrees of freedom permits unlimited occupation of a single state. Such a universe generically collapses into highly condensed configurations, lacking the rigidity and incompressibility observed in ordinary matter. In contrast, fermionic matter obeys the Pauli exclusion principle, which prevents multiple particles from occupying the same quantum state. This exclusion underlies electron shells, degeneracy pressure, and the solidity of everyday objects.

Any framework claiming that matter emerges from a more primitive quantum substrate must therefore explain how fermionic statistics arise dynamically. In the Hilbert Substrate Framework, the substrate consists of distinguishable qubits---no antisymmetry, particle labels, or spacetime structure are assumed \emph{a priori}. The central question is:

\begin{quote}
\emph{Can generic unitary dynamics on a qubit substrate produce effective quasiparticles that obey fermionic statistics?}
\end{quote}

\subsection{Fermionic Audit Metrics}

We introduce three independent diagnostics to determine whether fermionic matter has emerged.

\subsubsection{Sector Additivity (Fermi Sea Structure)}

We decompose the Hilbert space into magnetization sectors labeled by the number $n$ of spin-down excitations. These correspond to particle-number sectors in a fermionic interpretation.

Let $\varepsilon_k$ denote the eigenvalues of $H$ restricted to the one-particle sector ($n=1$). For free fermions, the energies in the two-particle sector ($n=2$) satisfy:
\begin{equation}
E^{(2)} \approx \varepsilon_{k_1} + \varepsilon_{k_2}, \quad k_1 \neq k_2.
\end{equation}

We compute the root-mean-square deviation between the actual two-particle energies and all allowed sums of distinct one-particle energies. Small residuals indicate that the many-body spectrum is well-approximated by a Fermi sea.

\subsubsection{Jordan--Wigner Anticommutation Test}

We construct fermionic operators from qubits using the Jordan--Wigner transformation:
\begin{equation}
c_j = \left( \prod_{m<j} Z_m \right) \sigma^-_j, \quad
c_j^\dagger = \left( \prod_{m<j} Z_m \right) \sigma^+_j.
\end{equation}

We evaluate expectation values of anticommutators on the recovered ground state $|\psi_0\rangle$:
\begin{align}
\langle \psi_0 | \{c_i, c_j\} | \psi_0 \rangle &\approx 0, \\
\langle \psi_0 | \{c_i, c_j^\dagger\} | \psi_0 \rangle &\approx \delta_{ij}.
\end{align}

\subsubsection{Pauli Pressure Proxy}

As a discrete analogue of degeneracy pressure, we examine the curvature of the ground-state energy as a function of particle number:
\begin{equation}
\kappa^{-1}(n) = E_0(n+1) - 2E_0(n) + E_0(n-1).
\end{equation}

For fermionic systems, $E_0(n)$ is convex due to the progressive filling of higher-energy modes. This convexity is absent in bosonic systems where particles can pile into the lowest mode.

\subsection{Numerical Results: Fermionic Emergence}

We study XX-model Hamiltonians on $N=8$ qubit chains:
\begin{equation}
H_{\text{XX}} = \sum_j \left( X_j X_{j+1} + Y_j Y_{j+1} \right).
\end{equation}

Our analysis does not rely on the known Jordan--Wigner equivalence. Instead, we:
\begin{enumerate}
\item Scramble the Hamiltonian by applying global unitary transformations.
\item Recover an accessible basis using geometry-blind optimization.
\item Audit the recovered Hamiltonian for fermionic behavior.
\end{enumerate}

\begin{result}[Emergent Fermionic Statistics]
Across multiple random scrambles and geometry-blind recoveries:
\begin{enumerate}
\item Jordan--Wigner anticommutation relations emerge with violations at or near machine precision ($\sim 10^{-14}$).
\item Many-body spectra exhibit additive structure characteristic of free fermions.
\item Ground-state energies display convexity consistent with Pauli exclusion.
\end{enumerate}
\end{result}

Crucially, these results hold even in runs where spatial locality is not cleanly recovered. This indicates that fermionic statistics are a more primitive emergent feature than geometry itself.

\subsection{Matter Before Geometry}

These findings establish an ordering of emergence:
\begin{equation}
\text{Hilbert substrate} \rightarrow \text{subsystems} \rightarrow \text{exchange statistics} \rightarrow \text{accessibility constraints} \rightarrow \text{locality} \rightarrow \text{geometry}.
\end{equation}

Fermionic diagnostics remain intact even in regimes where spatial locality is dynamically inaccessible. Global scrambling can destroy the path back to a spatially local basis, yet exchange statistics and sector structure persist. This demonstrates that matter-like structure can precede, and does not require, emergent spacetime locality.

%==============================================================================
\section{Constraint Separation and No-Refolding}
\label{sec:constraint_separation}
%==============================================================================

\subsection{Operational Definition of Matter}

To test the No-Refolding constraint, we require a precise operational definition of ``existing matter.'' We formalize this via a \emph{quadratic anchor}.

For a general Hamiltonian $H$, we extract its quadratic component by projecting onto the fermionic bilinear operator basis $\{c_i^\dagger c_j\}$:
\begin{equation}
t_{ij} = \frac{\text{Tr}\left[(c_i^\dagger c_j)^\dagger H\right]}{\text{Tr}\left[(c_i^\dagger c_j)^\dagger (c_i^\dagger c_j)\right]}.
\end{equation}

We reconstruct the best-fit quadratic Hamiltonian:
\begin{equation}
H_{\text{quad}} = \sum_{i,j} t_{ij}\, c_i^\dagger c_j.
\end{equation}

The \emph{quadratic fraction} is then:
\begin{equation}
Q(H) = 1 - \frac{\|H - H_{\text{quad}}\|}{\|H\|},
\end{equation}
where $\|\cdot\|$ denotes the Frobenius norm. By construction, $Q(H) = 1$ if and only if $H$ is exactly quadratic (free fermions).

\subsection{Locality Relative to Target Geometries}

Given a target graph $B$ on the same $N$ sites, we define a local operator subspace consisting of on-site Pauli operators and two-body Pauli products supported on the edges of $B$. The complementary nonlocal fraction is:
\begin{equation}
\text{leak}_B(H) = 1 - \frac{\|P_B(H)\|^2}{\|H\|^2},
\end{equation}
where $P_B$ projects onto the local subspace.

\subsection{The Constraint Separation Experiment}

Starting from a free-fermion hopping Hamiltonian $H_0$ on a ring graph $A$:
\begin{equation}
H_0 = -t \sum_{\langle i,j\rangle \in A} \left( c_i^\dagger c_j + c_j^\dagger c_i \right),
\end{equation}
we attempt to reduce $\text{leak}_B(H)$ (locality cost relative to a different target graph $B$) using only local unitary dynamics: conjugation by adjacent two-qubit gates.

We compare two regimes:

\textbf{Free refolding}: Minimize $\text{leak}_B(H)$ without constraints.

\textbf{No-refolding enforced}: Minimize $\text{leak}_B(H)$ subject to $Q(H) \geq 0.98$ at every step.

\begin{result}[Constraint Separation]
For target graphs disjoint from the source ring:
\begin{enumerate}
\item In the \emph{free} regime, $\text{leak}_B(H)$ can be reduced substantially (e.g., from 1.0 to $\sim$0.79) using only local gates. However, this reduction is accompanied by catastrophic loss of quadratic structure: $Q(H) \sim 10^{-3}$.
\item In the \emph{no-refolding} regime, every proposed local update violates the quadratic constraint and is rejected. The dynamics stalls entirely: $\text{leak}_B(H)$ remains at its initial value while $Q(H) = 1$ is preserved exactly.
\end{enumerate}
\end{result}

Both regimes satisfy no-signaling identically; the sole difference is whether existing matter structure must be preserved.

\subsection{Implication: No-Refolding as an Independent Constraint}

These observations establish:

\begin{quote}
There exist directions in Hilbert space that reduce signaling cost relative to a new geometry but are dynamically accessible under local evolution \emph{only if} existing fermionic structure is destroyed. When quadratic structure is preserved, those geometries become dynamically inaccessible.
\end{quote}

No-Refolding is therefore an independent constraint on Hilbert-space dynamics, distinct from No-Signaling. Once stable matter exists, the set of geometries reachable by local unitary evolution is strictly restricted.

Not all unitarily equivalent factorizations of Hilbert space are physically accessible once matter is present. Geometry becomes history-dependent: the presence of matter constrains which refactorizations can be dynamically realized.

%==============================================================================
\section{Dimensional Selection and the Settling Hypothesis}
%==============================================================================

\subsection{Accessibility versus Optimization}

As shown in Paper II, locality emerges because it is the simplest \emph{accessible} description of large quantum systems, not because it is globally optimal. The unitary orbit contains many equivalent descriptions, but only a subset are dynamically reachable under locality-respecting transformations.

We extend this insight to dimensionality. The accessibility landscape contains multiple geometric basins corresponding to different effective dimensions. The substrate does not survey these basins globally; instead, it \emph{settles} into the first basin capable of supporting persistent structure.

\begin{definition}[Settling Hypothesis]
Three-dimensional space exists not because it is globally optimal, but because it is the first dimension where stable fermionic matter can form. Once matter exists, No-Refolding prevents further geometric evolution. Dimensionality freezes at the first self-anchoring configuration.
\end{definition}

\subsection{Dimensional Requirements for Stable Matter}

Different spatial dimensions exhibit qualitatively distinct properties relevant to the stability of matter:

\textbf{One dimension.} Exchange statistics are nonlocal (Jordan--Wigner strings extend across the system), and stable bound structures are limited. The fundamental group of the configuration space is trivial.

\textbf{Two dimensions.} The braid group $B_n$ permits infinitely many exchange classes, allowing anyonic phases. While this permits richer statistics, it implies no intrinsic collapse to a minimal set of exchange sectors. Fermions in two dimensions are not self-protecting without additional topological order or energy gaps.

\textbf{Three dimensions.} The relevant configuration space of identical excitations has fundamental group $\pi_1 \cong \mathbb{Z}_2$. Only two homotopy classes of exchanges exist. There is no invariant capable of protecting a continuous $U(1)$ exchange phase. Fermionic exchange ($\phi = \pi$) is dynamically enforced as a stable fixed point. Coulombic bound states are stable, and topologically protected structures (knots) exist.

\textbf{Four or more dimensions.} Exchange topology already collapses to $\mathbb{Z}_2$, so higher dimensions do not provide additional protection for fermionic statistics. However, increased coordination accelerates nonlocal mixing, spinor representations become larger, and locality basins are empirically shallower.

\begin{proposition}
Three dimensions constitute the minimum spatial dimension supporting proper fermionic statistics, stable atomic bound states, and topologically persistent information storage.
\end{proposition}

\subsection{Exchange Statistics as a Dynamical Property}

In this framework, particle statistics are not imposed by symmetrization postulates. They are defined \emph{operationally} via the stability of exchange phases under local unitary evolution.

Let $\phi$ denote the phase accumulated under an operational exchange of two excitations. The stability functional $S(\phi)$ measures the persistence of this phase under admissible local dynamics:
\begin{equation}
S(\phi) \to 0 \quad \text{for} \quad \phi \neq 0, \pi.
\end{equation}

Only bosonic ($\phi = 0$) and fermionic ($\phi = \pi$) phases remain stable as system size grows. The mechanism is topological: in three dimensions, exchange loops can be continuously deformed and unwound, leaving only $\mathbb{Z}_2$-valued invariants.

%==============================================================================
\section{Numerical Evidence for Dimensional Selection}
%==============================================================================

\subsection{Experimental Design}

We perform locality-recovery experiments on Hamiltonians defined over interaction graphs of varying intrinsic dimension. Each experiment consists of:

\begin{enumerate}
\item Constructing a local XX-model Hamiltonian on a specified graph topology.
\item Applying a global scrambling unitary (Haar-random or deep random circuit).
\item Running a locality-biased recovery flow using graph-local two-qubit gates.
\item Measuring the reduction in nonlocal operator weight (locality cost).
\end{enumerate}

We compare: 1D chains, 2D square lattices, 3D cubic lattices, 4D hypercubic lattices, and degree-matched random regular graphs (non-geometric control).

\subsection{Results: Basin Depth by Dimension}

Figure~\ref{fig:dimensional_recovery} presents the central empirical result of this section. Across five random seeds per topology at $N=8$ qubits, we measure the improvement in local fraction achieved by the locality-biased recovery flow. The results demonstrate a clear trend: geometric dimensionality correlates with recovery performance.

\begin{figure}[H]
\centering
\includegraphics[width=0.85\textwidth]{fig_recovery_local_improvement_vs_graph.png}
\caption{\textbf{Locality recovery improves with geometric dimensionality.} Each point represents the change in local fraction ($\Delta$ local fraction) achieved by locality-biased recovery starting from a scrambled Hamiltonian, for $N=8$ qubits across 5 random seeds per graph topology. Error bars show mean $\pm$ standard deviation. Three-dimensional lattices exhibit the highest mean recovery and the largest variance, indicating deep but structured basins. Random regular graphs (degree-matched control) show intermediate recovery with high consistency, confirming that geometric structure---not merely connectivity---drives the dimensional advantage.}
\label{fig:dimensional_recovery}
\end{figure}

The 1D chain and 1D ring show modest, consistent recovery ($\Delta \approx 0.0003$). The 2D lattice improves substantially ($\Delta \approx 0.001$), while the 3D lattice achieves the highest recovery ($\Delta \approx 0.0014$ mean) with notably high variance across seeds. This variance is significant: it indicates that 3D geometry contains both deep basins (high-recovery runs) and shallower regions, consistent with a complex but ultimately favorable accessibility landscape.

The random regular graph (degree $d=4$, matched to 2D lattice coordination) serves as a non-geometric control. Its intermediate performance confirms that the dimensional advantage is not simply a consequence of increased connectivity, but reflects genuine geometric structure.

\begin{table}[h]
\centering
\begin{tabular}{lcccc}
\toprule
Topology & Dimension & Mean $\Delta$ Local & Std Dev & Basin Character \\
\midrule
1D chain & 1 & 0.00029 & 0.00013 & Shallow, consistent \\
1D ring & 1 & 0.00036 & 0.00022 & Shallow, variable \\
2D lattice & 2 & 0.00089 & 0.00042 & Moderate depth \\
3D lattice & 3 & 0.00140 & 0.00080 & Deep, structured \\
Random regular ($d{=}4$) & -- & 0.00072 & 0.00025 & Moderate, consistent \\
\bottomrule
\end{tabular}
\caption{Summary statistics for locality recovery across graph topologies ($N=8$, 5 seeds). Three-dimensional lattices exhibit both the highest mean recovery and the largest variance, indicating deep but complex accessibility basins.}
\label{tab:dimensional}
\end{table}

These results (Figure~\ref{fig:dimensional_recovery} and Table~\ref{tab:dimensional}) support the settling hypothesis: three-dimensional lattices provide the deepest accessible locality basins among the tested topologies. While higher-dimensional graphs are mathematically admissible, the 3D advantage is already clear at modest system sizes. The high variance in 3D recovery suggests a rich basin structure rather than a simple, featureless attractor.

\subsection{The Accessibility Collapse Revisited}

From Paper II, the accessibility collapse at $N \geq 5$ for 1D chains prevents access to the Harmonion (eigenbasis) minimum. Extending this analysis to higher-dimensional graphs reveals:

\begin{enumerate}
\item All geometric topologies exhibit accessibility collapse at sufficiently large $N$.
\item The critical $N$ at which collapse occurs depends on dimension.
\item 3D lattices show collapse to the \emph{deepest} local minimum among tested topologies.
\item Random (non-geometric) graphs show collapse to shallow, unstable basins.
\end{enumerate}

These results do not rely on spectral dimension, which is poorly resolved at small system sizes. Instead, they reflect structural accessibility: how readily a given geometry traps the locality recovery flow.

%==============================================================================
\section{Dimensional Lock-In via No-Refolding}
%==============================================================================

\subsection{The Settling Process}

Combining the results on fermionic emergence, constraint separation, and dimensional basins, we can now describe the complete settling process:

\begin{enumerate}
\item \textbf{Pre-geometric phase}: The substrate begins in a configuration with no preferred geometry. Information propagates according to the interaction topology, but no stable structures exist.

\item \textbf{Fermionic emergence}: Under accessibility-constrained dynamics, fermionic operators emerge as stable algebraic structures. Exchange statistics collapse to $\mathbb{Z}_2$ (bosons and fermions).

\item \textbf{Dimensional settling}: The substrate settles into a geometric basin. Three dimensions are selected because this is the first dimension where fermionic matter is both topologically protected and capable of anchoring geometry.

\item \textbf{No-Refolding lock-in}: Once matter exists, transitions to alternative geometries would require destroying the fermionic structure. Such transitions are dynamically forbidden.
\end{enumerate}

\subsection{Why Three Dimensions Are ``First''}

The phrase ``first geometry'' requires clarification. We do not mean temporally first in a cosmological sense, but rather \emph{first in the accessibility ordering}. As the substrate explores the unitary orbit under locality-biased dynamics, three dimensions constitute the first basin that:

\begin{enumerate}
\item Is deep enough to trap the dynamics (accessibility collapse).
\item Supports stable fermionic statistics (topological protection).
\item Allows fermionic matter to anchor geometry (No-Refolding activation).
\end{enumerate}

Lower dimensions fail condition (2) or (3). Higher dimensions fail to be reached before three-dimensional matter forms and activates lock-in.

\subsection{Hysteresis and History-Dependence}

The dimensional selection mechanism exhibits hysteresis:

\begin{enumerate}
\item \textbf{Before fermionic stabilization}: Effective dimensionality may fluctuate as the substrate explores geometric basins.
\item \textbf{After fermions emerge}: Dimensionality becomes locked. The system cannot transition to 2D or 4D without destroying the matter that anchors 3D geometry.
\end{enumerate}

This history-dependence explains why the question ``Why not 4D?'' has a different answer from ``Why not 2D?'': The substrate never reaches 4D because 3D matter forms first and prevents further transitions. Two dimensions fail to support the stable fermionic structures needed to activate lock-in.

%==============================================================================
\section{Discussion}
%==============================================================================

\subsection{Relation to Other Approaches}

This framework differs fundamentally from existing explanations of dimensionality:

\textbf{Anthropic arguments} note that stable atoms require three dimensions but do not explain how such dimensionality is selected. Our framework provides the dynamical mechanism.

\textbf{Compactification scenarios} in string theory presuppose a higher-dimensional geometric background. We require no such background; geometry emerges rather than being assumed.

\textbf{Holographic approaches} derive bulk geometry from boundary conformal field theory. While conceptually compatible, our framework does not require the specific mathematical structure of AdS/CFT.

\textbf{Emergent spacetime programs} (e.g., causal sets, spin foams) typically impose constraints designed to produce spacetime-like structure. Our constraints (No-Signaling, No-Forgetting, No-Refolding) are more minimal and information-theoretic.

The distinctive feature of our approach is that dimensionality is selected by \emph{accessibility} rather than optimality, energetics, or observer selection. Higher-dimensional geometries are not forbidden in principle; they are simply never reached.

\subsection{The Complete Emergence Chain}

We can now state the complete emergence chain within the Hilbert Substrate Framework:

\begin{enumerate}
\item \textbf{Hilbert space} $\mathcal{H}$ with unitary dynamics (fundamental).
\item \textbf{Subsystem factorization} (operational, via interaction graphs).
\item \textbf{Heterogeneous information propagation} (Paper I: HIP persistence).
\item \textbf{Accessibility collapse} (Paper II: phase transition to locality).
\item \textbf{Fermionic statistics} (this paper: emergent anticommutation).
\item \textbf{Dimensional selection} (this paper: 3D as first stable basin).
\item \textbf{No-Refolding lock-in} (this paper: matter anchors geometry).
\end{enumerate}

Each step is derived from the previous rather than assumed. The endpoint is three-dimensional space containing fermionic matter with local interactions---precisely the structure observed in physical reality.

\subsection{Physical Reality as Compromise}

The central philosophical conclusion of this work is that physical reality represents a \emph{compromise} rather than an optimum.

The globally optimal description of an isolated quantum system is its energy eigenbasis (the ``Harmonion'' vacuum), where all degrees of freedom are independent and non-interacting. Our universe does not exhibit this structure because:

\begin{enumerate}
\item For large systems, the eigenbasis is kinetically inaccessible (Paper II).
\item The accessible basins contain interacting structure.
\item Three-dimensional geometry provides the deepest accessible basin.
\item Once matter forms in this basin, lock-in prevents further optimization.
\end{enumerate}

Physical reality is therefore not the simplest possible description of the universe, but the simplest description accessible to systems containing stable matter. Space, locality, and three-dimensionality are not fundamental but emergent---``good enough'' solutions that the universe settles into because perfect solutions are out of reach.

\subsection{Limitations and Future Directions}

Several limitations of the present work should be noted:

\textbf{System size}: Numerical experiments are limited to $N \leq 14$ qubits. Verifying scaling behavior toward thermodynamic limits will require tensor-network methods or analytical bounds.

\textbf{Continuous dimensions}: We have compared discrete dimensional lattices. Understanding the transition between integer dimensions remains open.

\textbf{Cosmological dynamics}: We have not addressed how the settling process relates to cosmological evolution or whether signatures of pre-geometric phases might be observable.

\textbf{Forces and particles}: While we have demonstrated emergent locality and fermionic statistics, deriving the specific content of the Standard Model (gauge symmetries, particle masses, coupling constants) remains a challenge for future work.

Future directions include:
\begin{itemize}
\item Analytical bounds on basin depth as a function of dimension.
\item Tensor-network simulations of larger substrates.
\item Investigation of gauge structure emergence from accessibility constraints.
\item Cosmological implications of dimensional settling and potential observational signatures.
\end{itemize}

%==============================================================================
\section{Conclusion}
%==============================================================================

We have demonstrated that three-dimensional space, fermionic matter, and local interactions emerge naturally within the Hilbert Substrate Framework from unitary dynamics subject to minimal information-theoretic constraints.

The key results are:

\begin{enumerate}
\item \textbf{Emergent fermionic statistics}: Jordan--Wigner operators arise dynamically on a qubit substrate, satisfying anticommutation relations without any antisymmetrization postulate. Fermionic matter (``steel'') emerges from a substrate containing no fundamental fermionic degrees of freedom.

\item \textbf{No-Refolding as a constraint}: Once stable fermionic structure exists, geometric transitions that would destroy this structure become dynamically inaccessible. This constraint is distinct from No-Signaling and emerges from applying information-theoretic consistency to substrates containing matter.

\item \textbf{Dimensional selection via settling}: Three dimensions are selected not by optimization but by accessibility. The substrate settles at the first geometric basin capable of supporting stable fermionic matter. Once such matter exists, No-Refolding locks the geometry in place.
\end{enumerate}

The universe does not choose three dimensions because they are optimal. It remains three-dimensional because this is where the settling process ends.

Physical reality is a compromise between the abstract simplicity of Hilbert space and the constraints required for observation. Space is not a fundamental background but an emergent attractor---a ``glassy'' state of the Hilbert substrate that arises because the universe is too complex to find its own simplest basis, and too constrained by its own matter to continue searching.

%==============================================================================
\section*{Acknowledgments}
%==============================================================================

[To be added.]

%==============================================================================
\bibliographystyle{plain}
\begin{thebibliography}{99}

\bibitem{HSF1}
B.~Bray.
\newblock The Hilbert Substrate Framework: Persistent Heterogeneity in Information Propagation.
\newblock \emph{Zenodo}, 2025.
\newblock DOI: 10.5281/zenodo.18076311.

\bibitem{HSF2}
B.~Bray.
\newblock The Hilbert Substrate Framework II: Accessibility and the Emergence of Spatial Locality.
\newblock \emph{Zenodo}, 2025.

\bibitem{Nielsen}
M.~A.~Nielsen and I.~L.~Chuang.
\newblock \emph{Quantum Computation and Quantum Information}.
\newblock Cambridge University Press, 10th anniversary edition, 2010.

\bibitem{Watrous}
J.~Watrous.
\newblock \emph{The Theory of Quantum Information}.
\newblock Cambridge University Press, 2018.

\bibitem{LiebRobinson}
E.~H.~Lieb and D.~W.~Robinson.
\newblock The finite group velocity of quantum spin systems.
\newblock \emph{Communications in Mathematical Physics}, 28:251--257, 1972.

\bibitem{Brockett}
R.~W.~Brockett.
\newblock Dynamical systems that sort lists, diagonalize matrices, and solve linear programming problems.
\newblock \emph{Linear Algebra and its Applications}, 146:79--91, 1991.

\bibitem{Wales}
D.~J.~Wales.
\newblock \emph{Energy Landscapes}.
\newblock Cambridge University Press, 2003.

\bibitem{Wen}
X.-G.~Wen.
\newblock Choreographed entanglement dances: Topological states of quantum matter.
\newblock \emph{Science}, 363(6429):eaal3099, 2019.

\bibitem{Zurek}
W.~H.~Zurek.
\newblock Quantum Darwinism.
\newblock \emph{Nature Physics}, 5:181--188, 2009.

\bibitem{Bravyi}
S.~Bravyi, M.~B.~Hastings, and F.~Verstraete.
\newblock Lieb-Robinson bounds and the generation of correlations and topological quantum order.
\newblock \emph{Physical Review Letters}, 97:050401, 2006.

\end{thebibliography}

\end{document}